% =============================================================================
%  I. INTRODUCTION: THE CLINICAL PROBLEM
% =============================================================================
\begin{frame}{I. Introduction: The Clinical Problem}
\begin{itemize}
  \item \textbf{Sepsis} is a life-threatening condition caused by the body's
        extreme response to an infection~\cite{singer2016sepsis3}.
  \item It causes \textbf{about 11 million deaths every year}, nearly
        \textbf{1 in 5 deaths worldwide}~\cite{rudd2020global}.
  \item Delaying treatment by even a few hours increases the risk of
        death~\cite{seymour2017time}.
  \item Early warning systems aim to detect sepsis before doctors recognise
        it, allowing earlier treatment.
\end{itemize}

\vspace{8pt}
\textbf{More than 90 AI models} have already been developed to predict
sepsis~\cite{wang2025srpm}.

\vfill
\begin{block}{}
  \centering Existing models have improved prediction, but important
  methodological questions remain unanswered.
\end{block}
\end{frame}


% =============================================================================
%  I. INTRODUCTION: RESEARCH QUESTION
% =============================================================================
\begin{frame}{I. Introduction: Research Question}
\vfill
\begin{block}{RQ}
  In real-time sepsis onset prediction on MIMIC-IV, how much of the observed
  performance variation is explained by the choice of model, the sepsis label
  definition, and the analyst's methodological decisions?
\end{block}
\vfill
\end{frame}
